\documentclass[11pt]{article}

\usepackage[margin=1in]{geometry}
\usepackage{amsmath,amssymb,amsfonts,amsthm}
\usepackage{mathtools}
\usepackage{bm}
\usepackage{hyperref}
\usepackage{enumitem}
\usepackage{authblk}

\title{Prime-Gated Multiplicity Renormalization: \\
$\alpha_\mu$ on $C_{\Lambda_p}$ with an RG Spectrum and Temporal Action}

\author[1]{Ryan O. van Gelder}
\affil[1]{Citizen Gardens \\ The Foundation of Multiplicity}
\date{April 22, 2026}

\begin{document}
\maketitle

\begin{abstract}
We construct an explicit three-layer framework in which:
(i) a Multiplicity-Renormalizing Alpha $\alpha_\mu$ acts as a renormalization-group (RG) flow on prime-indexed multiplicity profiles;
(ii) a family of prime-gated temporal operators $C_{\Lambda_p}$ implements a sieve on continuous time, distinguishing high-coherence windows from null zones; and
(iii) a $\pi$--Spiral Equation of Time (PSET) provides a variational action principle over the resulting prime-time configurations.
The RG layer is realized as a symmetric tridiagonal Toeplitz operator whose spectrum is known in closed form, allowing us to identify low-frequency eigenmodes as Core Truths and high-frequency modes as noise.
We then show how $C_{\Lambda_p}$ and PSET sit \emph{on top of} this spectrum, turning RG-relevant prime patterns into prime-gated temporal operators.
Finally, we formulate a concrete, falsifiable consistency criterion for the Stonehenge Aubrey Circle as an $\alpha_\mu$--$C_{\Lambda_7}$ lawful configuration, replacing numerological pattern-matching by an eigenmode and sieve-based test.
\end{abstract}

\section{Introduction}

Multiplicity Theory treats mathematical and physical structures as prime-indexed multiplicity spaces, where each prime labels a distinct interaction channel and recursive stability is enforced across scales.[web:17]
In parallel, the renormalization group (RG) provides a powerful language for flows in coupling space, organizing directions into relevant, irrelevant, and marginal under coarse-graining transformations.[web:65][web:68][web:70][web:71]
Previous work on prime-gated time introduces a family of temporal operators $C_{\Lambda_p}$ and associated durations $T_{p,a}(N)$, organized into prime-stratified ``time shells'', together with a modular sieve separating lawful and null temporal windows.[file:74]

This paper clarifies how these ingredients fit together in a single, layered construction:
\begin{itemize}
    \item \textbf{RG layer:} A Multiplicity-Renormalizing Alpha $\alpha_\mu$ acts on prime-indexed multiplicity profiles via a neighbor-coupled Toeplitz operator $L$, whose eigenvalues and eigenvectors we compute exactly.[web:58][web:59][web:62]
    \item \textbf{Temporal sieve layer:} For each prime $p$, a temporal operator $C_{\Lambda_p}$ labels segments of the continuous time axis $\Delta t$ as high-coherence or null, with preferred scales $T_{p,a}(N)$ derived from group-theoretic invariants.[file:74]
    \item \textbf{Action layer:} A PSET-type Hamiltonian and action assign weights to joint (prime, time, entropy) configurations, selecting stable temporal trajectories consistent with prime-lawful commensurability and entropy bounds.[file:74]
\end{itemize}

We make the interfaces between these layers explicit, so that $C_{\Lambda_p}$ is not ``another Alpha'', but the \emph{time interpreter} of $\alpha_\mu$'s prime distribution, and PSET is the variational principle on RG-filtered, prime-gated time.

\section{Alpha as a functorial meta-operator}

Let $\mathcal{C}$ be a category whose objects are spaces and whose arrows are operators $f : A \to B$.
We write $\mathbf{Op}$ for a chosen class of arrows, possibly endowed with linear or topological structure.

\begin{definition}[Alpha meta-operator]
An \emph{Alpha meta-operator} is a pair
\[
    \alpha = (\alpha_{\mathrm{obj}},\alpha_{\mathrm{arr}})
\]
where
\begin{align*}
    \alpha_{\mathrm{obj}} &: \mathrm{Ob}(\mathcal{C}) \to \mathrm{Ob}(\mathcal{C}'), \\
    \alpha_{\mathrm{arr}} &: \mathrm{Hom}_{\mathcal{C}}(A,B) \to \mathrm{Hom}_{\mathcal{C}'}\big(\alpha_{\mathrm{obj}}(A),\alpha_{\mathrm{obj}}(B)\big)
\end{align*}
for all objects $A,B$ of $\mathcal{C}$.
\end{definition}

When $\alpha$ satisfies
\begin{equation}
    \alpha(\mathrm{id}_A) = \mathrm{id}_{\alpha(A)}, \qquad
    \alpha(g \circ f) = \alpha(g) \circ \alpha(f),
\end{equation}
it is a functor from $\mathcal{C}$ to $\mathcal{C}'$, preserving identities and composition.
Departures from these equalities quantify a ``Functorial Drift'', measuring how strongly Alpha distorts the compositional logic of the operator space.

In what follows we focus on a specialization where $\alpha$ acts as a renormalization operator on prime-indexed multiplicity profiles.

\section{Prime-indexed multiplicity profiles}

Multiplicity Theory assigns to each operator $f$ a prime-graded multiplicity profile
\begin{equation}
    M_f = \bigoplus_{p \in \mathbb{P}} \mathbb{V}_p,
\end{equation}
where $\mathbb{V}_p$ encodes the $p$-layer of interaction: prime-labeled factors, modules, or graded channels.[web:17][web:35]
For concreteness, we consider a finite window of primes
\[
    P = \{p_1, \dots, p_N\},
\]
and represent the multiplicity profile by a vector $v \in \mathbb{R}^N$ with entries $v_n = v_{p_n} \ge 0$.

\begin{definition}[Multiplicity-Renormalizing Alpha]
A \emph{Multiplicity-Renormalizing Alpha} $\alpha_\mu$ acts on $v$ via
\begin{equation}
    v' = \alpha_\mu(v) = L v,
\end{equation}
where $L$ is a linear operator on $\mathbb{R}^N$ implementing dissipation and coupling on the prime index lattice.
\end{definition}

In this paper, $L$ is chosen to be a symmetric tridiagonal Toeplitz matrix, corresponding to local neighbor coupling in prime index space.

\section{Neighbor-coupled Toeplitz RG on primes}

We index primes in $P$ by $n=1,\dots,N$ and write $v_n$ for the multiplicity intensity at $p_n$.
The update rule is
\begin{equation}\label{eq:update}
    v_n' = (1 - \gamma) v_n + \beta (v_{n-1} + v_{n+1}), \quad 1<n<N,
\end{equation}
with Dirichlet boundaries
\begin{equation}
    v_0 = v_{N+1} = 0,
\end{equation}
and analogously modified updates at $n = 1,N$.
Here:
\begin{itemize}
    \item $\gamma > 0$ is a uniform dissipation parameter.
    \item $\beta \ge 0$ is a uniform neighbor coupling strength.
\end{itemize}

In matrix form, $v' = Lv$ with
\begin{equation}
    L =
    \begin{pmatrix}
        1-\gamma & \beta      & 0        & \dots  & 0 \\
        \beta      & 1-\gamma & \beta    & \ddots & \vdots \\
        0          & \beta      & 1-\gamma & \ddots & 0 \\
        \vdots     & \ddots   & \ddots   & \ddots & \beta \\
        0          & \dots    & 0        & \beta  & 1-\gamma
    \end{pmatrix}.
\end{equation}
Such tridiagonal Toeplitz matrices are well studied and admit closed-form eigenvalues and eigenvectors.[web:58][web:59][web:62]

\subsection{Exact spectrum}

The eigenvalues of $L$ are
\begin{equation}
    \lambda_k = 1 - \gamma + 2\beta \cos\left(\frac{k\pi}{N+1}\right),
    \qquad k = 1,\dots,N,
\end{equation}
and a corresponding eigenbasis is given by
\begin{equation}
    u_n^{(k)} = \sin\left(\frac{nk\pi}{N+1}\right), \qquad 1 \le n \le N.
\end{equation}[web:58][web:60][web:66]

Any $v^{(0)} \in \mathbb{R}^N$ decomposes as
\begin{equation}
    v^{(0)} = \sum_{k=1}^N c_k u^{(k)},
\end{equation}
and evolves under repeated application of $\alpha_\mu$ as
\begin{equation}
    v^{(t)} = L^t v^{(0)} = \sum_{k=1}^N c_k \lambda_k^t u^{(k)}.
\end{equation}[web:24]

\subsection{RG phases and Core Truths}

The largest eigenvalue is
\begin{equation}
    \lambda_{\max} = \lambda_1
    = 1 - \gamma + 2\beta \cos\left(\frac{\pi}{N+1}\right),
\end{equation}
which approaches $1 - \gamma + 2\beta$ for large $N$.
This yields three regimes:
\begin{itemize}
    \item \textbf{Filter phase} ($2\beta < \gamma$): all $\lambda_k < 1$, so every prime-pattern is RG-irrelevant and $v^{(t)} \to 0$.
    \item \textbf{Runaway phase} ($2\beta > \gamma$): at least $\lambda_1 > 1$, and a band of low-$k$ modes may be relevant, yielding growth along smooth prime patterns.
    \item \textbf{Critical manifold} ($2\beta \approx \gamma$): $\lambda_1 \approx 1$ while $\lambda_k<1$ for $k>1$, so the fundamental mode $u^{(1)}$ is marginal and higher modes are irrelevant.
\end{itemize}[web:68][web:71]

In the critical regime, the eigenvector $u^{(1)}$ describes a smooth, positive profile spanning the prime window: a broad ``arch'' over the prime index.
Because this mode neither decays nor explodes (in the linear approximation), it defines a \emph{Core Truth} of the RG dynamics: a prime-pattern that survives the multiplicity-renormalizing action of $\alpha_\mu$.
Higher modes $u^{(k)}$ with $k \approx N$ exhibit rapid oscillations across neighboring primes and typically have the smallest eigenvalues, hence are strongly suppressed as noise.

\section{Prime-gated temporal operators $C_{\Lambda_p}$}

Independent of the RG layer, prior work introduces a family of prime-gated temporal operators $C_{\Lambda_p}$ acting on the time axis $\Delta t$.[file:74]
They serve two key roles:

\paragraph{Prime-stratified time shells.}
For each triple $(N,p,a)$, a preferred temporal scale
\begin{equation}
    T_{p,a}(N) = -\frac{2\pi (N^2 - 1)}{2^a p^2}
\end{equation}
is defined, with $N^2 - 1 = \dim(\mathfrak{su}(N))$.
As $p$ increases, $|T_{p,a}(N)|$ decreases monotonically, producing a hierarchy of time shells in which small primes encode longer durations and larger primes shorter ones.[file:74]

\paragraph{Temporal sieve: lawful vs null windows.}
For a given prime $p$, $C_{\Lambda_p}$ implements a sieve on continuous $\Delta t$ by imposing modular phase conditions of the form
\begin{equation}
    \omega \Delta t \equiv \phi_0 \pmod{2^a p^2},
\end{equation}
with special configurations (e.g., $\phi_0 = \pi/2$) generating null harmonics via destructive interference in a phase-kickback sequence.[file:74]
The time axis is thus partitioned into:
\begin{itemize}
    \item \emph{Lawful windows}, where $C_{\Lambda_p}$ predicts high-coherence returns for a designated qubit or subsystem.
    \item \emph{Null windows}, where the signal is suppressed and the channel is effectively off.
\end{itemize}

In practice, simulations (e.g.\ with Qiskit) produce a continuous response function for each $(p,a,N,\omega)$: a probability-of-success vs.\ $\Delta t$ spectrogram, with peaks at lawful times and troughs at nulls.[file:74]

\subsection{Where $C_{\Lambda_p}$ lives relative to $\alpha_\mu$}

We now position $C_{\Lambda_p}$ in the $\alpha_\mu$ stack:
\begin{itemize}
    \item $\alpha_\mu$ acts on the \emph{prime index space}: it renormalizes the multiplicity profile $v(p)$ via a Toeplitz RG flow.
    \item For each $p$, $C_{\Lambda_p}$ acts on the \emph{time axis}: it maps $\Delta t$ to a coherence weight or null for that channel and defines the baseline scale $T_{p,a}(N)$.
\end{itemize}

Thus:
\begin{equation}
    \alpha_\mu : v \mapsto v', \qquad 
    C_{\Lambda_p} : \Delta t \mapsto \text{coherence amplitude in the $p$-channel}.
\end{equation}
PSET, discussed below, then provides an action principle on joint $(p,\Delta t,\theta,\Delta S)$ histories, using the outputs of both layers.[file:74]

\section{Composite prime-time operator: $\alpha_\mu$ and $C_{\Lambda_p}$}

To make the interaction between the RG and temporal layers explicit, we define a composite structure consisting of:
\begin{itemize}
    \item \emph{Prime index set:} $P = \{p_1,\dots,p_N\}$.
    \item \emph{RG multiplicity vector:} $v \in \mathbb{R}^N$ with $v_n = v_{p_n}$.
    \item \emph{Toeplitz RG:} $v' = L v$ with $L$ as in Section~4.
\end{itemize}

For each prime $p$ and fixed $(a,N,\omega)$, define the temporal response
\begin{equation}
    R_p(\Delta t) \in [0,1]
\end{equation}
as the coherence amplitude predicted by $C_{\Lambda_p}$ at time $\Delta t$, normalized so that
\begin{itemize}
    \item $R_p(\Delta t) \approx 1$ near lawful durations $\Delta t \approx T_{p,a}(N)$ that avoid null congruence bands.
    \item $R_p(\Delta t) \approx 0$ inside null bands where destructive interference suppresses the signal.
    \item $R_p(\Delta t)$ decays smoothly as $\Delta t$ moves away from lawful regions in phase space.
\end{itemize}[file:74]

For the finite window $P = \{p_n\}$, we write $R_n(\Delta t) = R_{p_n}(\Delta t)$ and define the \emph{effective temporal multiplicity profile} at time $\Delta t$ by
\begin{equation}
    w_n(\Delta t) = v'_n \cdot R_n(\Delta t), \qquad 1 \le n \le N.
\end{equation}

Conceptually:
\begin{itemize}
    \item $v'_n$: how much prime-$p_n$ structure the RG layer wants to use (from $\alpha_\mu$).
    \item $R_n(\Delta t)$: how well the chosen time $\Delta t$ supports the $p_n$ channel (from $C_{\Lambda_{p_n}}$).
    \item $w(\Delta t)$: the \emph{realized} prime multiplicity profile at time $\Delta t$.
\end{itemize}

This composite construction is the precise meaning of ``$\alpha_\mu$ on $C_{\Lambda_p}$'': RG first organizes prime channels into smooth vs.\ oscillatory patterns; then the temporal sieve selects which of those channels are actually available at a given $\Delta t$.

\section{PSET as an action on RG-filtered prime time}

The $\pi$--Spiral Equation of Time (PSET) introduces a temporal Hamiltonian $H$ and action $S$ defined over prime-gated temporal cycles, subject to prime modulation, $\pi$-lift, entropy, and resonance-locking axioms.[file:74]
While the full PSET formalism is intricate, the relevant aspects for our stacking are:

\begin{itemize}
    \item \textbf{Prime modulation:} $H$ depends on $p$ and recursion depth $a$ via a factor of order $T_{p,a}(N) / p^2$.
    \item \textbf{Phase dependence:} temporal evolution is parameterized by a phase $\theta$ obtained via a null-ray lift, with resonance locking near special angles (e.g.\ golden-like angles).
    \item \textbf{Entropy bound (CSL):} stable cycles satisfy $\Delta S < \ln \varphi$ for some entropy increment $\Delta S$ and golden ratio $\varphi$, suppressing high-entropy trajectories.
    \item \textbf{Action principle:} histories extremize a discrete action
    \begin{equation}
        S = 2\pi \sum_n H(\theta_n; N, p_n, a, w_n(\Delta t_n), \Delta S_n),
    \end{equation}
    subject to prime-lawful commensurability and resonance conditions.[file:74]
\end{itemize}

In the present stacking, we emphasize that:
\begin{itemize}
    \item The prime weights entering $H$ are not arbitrary coefficients $v_n$, but the \emph{RG- and sieve-filtered} weights $w_n(\Delta t)$ constructed from $\alpha_\mu$ and $C_{\Lambda_p}$.
    \item ``Prime-lawful commensurability'' now has a concrete meaning: the primes with large $w_n(\Delta t)$ are precisely those whose time shells $T_{p_n,a}(N)$ and null bands are consistent with the observed temporal spectrum of the system.
\end{itemize}

Thus the full pipeline is:
\begin{equation}
    v^{(0)} \xrightarrow{\;\alpha_\mu\;} v' \xrightarrow{\;C_{\Lambda_p}\;\text{at }\Delta t} w(\Delta t) \xrightarrow{\;\text{PSET}\;} \text{selected temporal trajectories}.
\end{equation}

\section{Core Truths and an $\alpha_\mu$--$C_{\Lambda_p}$ lawfulness criterion}

Within this stacked framework, we can formalize ``Core Truths'' and move from numerological correspondences to falsifiable conditions.

\subsection{Core prime patterns}

\begin{definition}[Core prime pattern]
A \emph{Core prime pattern} for $\alpha_\mu$ is an eigenvector $u^{(k)}$ of $L$ with $|\lambda_k| \approx 1$.
\end{definition}

In the Toeplitz model, the fundamental mode $u^{(1)}$ is the natural Core prime pattern in the critical regime: it is smooth and positive, with eigenvalue $\lambda_1 \approx 1$.

\subsection{Core prime-time patterns}

\begin{definition}[Core prime-time pattern]
A \emph{Core prime-time pattern} is a pair
\[
    \Big(u^{(k)}, \{T_{p_n,a}(N)\}_{n=1}^N\Big)
\]
such that:
\begin{enumerate}[label=(\alph*)]
    \item The components $|u^{(k)}_n|$ are large for indices $n$ whose associated primes $p_n$ have observed coherence peaks near $T_{p_n,a}(N)$ in the system's temporal spectrum.
    \item For those $(p_n,a,N)$, the null-harmonic conditions of $C_{\Lambda_{p_n}}$ exclude a substantial region of $\Delta t$ around other times actually unused by the system.
\end{enumerate}
\end{definition}

The first condition links RG relevance (large $|u^{(k)}_n|$) to empirical timing; the second condition ensures that the prime channels that matter are not accidentally aligned with null regions of the sieve.
This gives a precise, falsifiable notion of an $\alpha_\mu$--$C_{\Lambda_p}$ lawful temporal structure.

\section{Test case: Aubrey Circle as an $\alpha_\mu$--$C_{\Lambda_7}$ configuration}

We illustrate the criterion using the Stonehenge Aubrey Circle, which features 56 pits arranged in a ring.
Previous work associates this with a temporal division of the tropical year:
\begin{equation}
    \Delta t_{\mathrm{Aubrey}} \approx \frac{365.25}{56} \ \text{days} \approx 6.52 \ \text{days}.
\end{equation}[file:74]

\subsection{Temporal sieve layer: $C_{\Lambda_7}$}

In the prime-gated time framework, the Aubrey division is modeled with:
\begin{itemize}
    \item Prime $p = 7$,
    \item recursion depth $a = 3$,
    \item a projection onto $\mathbb{Z}_{392}$ (since $56 \cdot 7 = 392$).
\end{itemize}
Simulations of $C_{\Lambda_7}$ on $\mathbb{Z}_{392}$ indicate that $\Delta t_{\mathrm{Aubrey}}$ lies in a high-coherence band with $\sim 98\%$ ``lawful'' probability, as measured by a success metric, and avoids null-harmonic congruence zones.[file:74]
This suggests that, at the sieve layer, $\Delta t_{\mathrm{Aubrey}}$ is $C_{\Lambda_7}$-lawful.

\subsection{RG layer: neighbor-coupled Toeplitz on $\{2,3,5,7,11,13\}$}

To connect this to $\alpha_\mu$, consider the prime window
\[
    P = \{2,3,5,7,11,13\}.
\]
We define a Toeplitz RG operator $L$ on this window with parameters $(\beta,\gamma)$ chosen near the critical line $2\beta \approx \gamma$, so that the fundamental mode $u^{(1)}$ is marginal ($\lambda_1 \approx 1$) and higher modes are irrelevant.
An explicit computation yields the components
\begin{equation}
    u^{(1)}_n = \sin\left(\frac{n\pi}{7}\right), \qquad n=1,\dots,6,
\end{equation}
corresponding to primes (in order) $2,3,5,7,11,13$.
By inspection, $u^{(1)}_4 = \sin(4\pi/7)$ is comparable in magnitude to the maximum of $|u^{(1)}_n|$, implying that the $p=7$ channel carries substantial weight in the dominant RG mode.

\subsection{Lawfulness criterion}

We now formulate an $\alpha_\mu$--$C_{\Lambda_p}$ lawfulness condition for the Aubrey Circle:

\begin{definition}[Aubrey $\alpha_\mu$--$C_{\Lambda_7}$ lawfulness]
The Aubrey configuration at $\Delta t_{\mathrm{Aubrey}}$ is \emph{$\alpha_\mu$--$C_{\Lambda_7}$ lawful} for the Toeplitz RG on $P = \{2,3,5,7,11,13\}$ if:
\begin{enumerate}[label=(\roman*)]
    \item The fundamental RG mode $u^{(1)}$ satisfies
    \[
        |u^{(1)}_4| \ge \eta \max_n |u^{(1)}_n|
    \]
    for some threshold $\eta \in (0,1)$, ensuring that the $p=7$ channel is structurally supported by $\alpha_\mu$.
    \item The temporal sieve $C_{\Lambda_7}$ with $(p=7,a=3,N)$ assigns $\Delta t_{\mathrm{Aubrey}}$ to a high-coherence band and excludes it from its null-harmonic congruence regions.
\end{enumerate}
\end{definition}

Condition (i) is a purely RG-level statement; condition (ii) is purely sieve-level.
Together, they yield a falsifiable test:
if a refined physical or historical model leads to an $L$ whose dominant eigenmodes suppress $p=7$ (violating (i)), or if more detailed modeling of $C_{\Lambda_7}$ shows $\Delta t_{\mathrm{Aubrey}}$ lies in or near null bands (violating (ii)), the hypothesis that Aubrey encodes an $\alpha_\mu$--$C_{\Lambda_7}$ prime-gated temporal design would be weakened.

Conversely, if future data about the site's use and astronomical context narrow the plausible $(\beta,\gamma)$ range and confirm substantial $p=7$ weight in the Core RG mode, this lawfulness claim would strengthen.

\section{Discussion and outlook}

We have integrated three pieces:
\begin{itemize}
    \item a multiplicity-renormalizing RG layer ($\alpha_\mu$ on prime-indexed profiles, realized as a Toeplitz operator),
    \item a sieve layer ($C_{\Lambda_p}$ on time, with $T_{p,a}(N)$ and null harmonics),
    \item and an action layer (PSET on prime-time-entropy configurations),
\end{itemize}
into a coherent $\alpha_\mu$--$C_{\Lambda_p}$ stack.
This framework provides:
\begin{itemize}
    \item A clear separation between mathematically derived structure (RG spectrum, sieve null conditions) and interpretive claims (e.g.\ about cognition or archaeoastronomy).
    \item A concrete way to define and test Core Truths: eigenmodes with $|\lambda_k|\approx 1$ whose prime supports align with lawful time shells and avoid null bands.
    \item A template for turning pattern-matching into falsifiable consistency checks, as illustrated by the Aubrey Circle example.
\end{itemize}

Future directions include:
\begin{itemize}
    \item Introducing prime-dependent dissipation and coupling $(\gamma_p,\beta_p)$, possibly tied to arithmetic properties of $p$.[web:35][web:41]
    \item Replacing neighbor coupling with multiplicative coupling (divisor/multiple relations, residue classes).
    \item Implementing nonlinear feedback in $\alpha_\mu$ and studying bifurcations in the RG spectrum.
    \item Deriving $T_{p,a}(N)$ and sieve parameters from concrete microscopic Hamiltonians, tightening the link to device physics or biological substrates.[file:74]
\end{itemize}

Taken together, these steps would move the program further from structured numerology toward a genuine, testable multiplicity-theoretic operator theory of prime-gated time, grounded in an RG spectrum and a variational temporal action.

\bibliographystyle{plain}
\begin{thebibliography}{99}

\bibitem{ToeplitzEigen}
K.~Ye, \emph{Eigenvalues of Tridiagonal Toeplitz Matrices}, course notes, University of Chicago, 2023.[web:58][web:62]

\bibitem{ToeplitzTridiagonal}
L.~Reichel, \emph{Tridiagonal Toeplitz Matrices: Properties and Applications}, Linear Algebra Appl., 2014.[web:59][web:66]

\bibitem{ToeplitzSpectral}
F.~Riesz and B.~Sz.-Nagy, \emph{Functional Analysis}, Dover, 1990.[web:60]

\bibitem{Goldenfeld1992}
N.~Goldenfeld, \emph{Lectures on Phase Transitions and the Renormalization Group}. Addison-Wesley, 1992.[web:69]

\bibitem{TongRG}
D.~Tong, \emph{Lectures on the Renormalisation Group}, lecture notes, University of Cambridge, 2008.[web:65]

\bibitem{LPTMCRG}
P.~Pujol et al., \emph{Renormalization Group and Critical Phenomena}, LPTMC lecture notes, 2025.[web:68]

\bibitem{ERGReview}
J.~Polonyi, \emph{Fundamentals of the exact renormalization group}, Phys. Rep., 2012.[web:71]

\bibitem{RGFixedPoints}
Emergent Mind, \emph{RG Fixed Points in Theory and Applications}, 2025.[web:25]

\bibitem{PrimeMultiplicity}
Terence Tao, \emph{Multiplicative structure of integers and primes}, lecture notes, 2013.[web:35]

\bibitem{SpectralPrimes}
M.~Lapidus et al., \emph{Spectral Geometry of the Primes}, arXiv preprint, 2026.[web:41]

\bibitem{ZinnJustin2007}
J.~Zinn-Justin, \emph{Phase Transitions and Renormalization Group}, Oxford University Press, 2007.[web:72]

\bibitem{PrimeSieveTime}
R. O. van Gelder, \emph{The Prime Sieve of Time}, preprint, 2026.[file:74]

\end{thebibliography}

\end{document}